\section{Introduction}

System administrators extend their software systems to customize it for their needs.  These extensions may tailor the application for hardware or a particular workload, as exemplified by works that improve system performance through extensions~\cite{Ghigoff21, Zhong22}.  Extensions may also enable system observability, as exemplified by works that extend systems to improve security~\cite{Jia23} or monitor performance~\cite{Gregg19, shen2023network,Zhu23}.  

For example, consider a system administrator who aims to understand the role of low-level library and kernel behavior on the tail latency of a microservice application.  The administrator could use system extensions to inject monitoring code across their software stack.  They would track requests by injecting monitoring logic into the application, track synchronization primitives by injecting logic into dependent libraries (\eg the \texttt{go} runtime, \texttt{libc}), and track kernel structures by injecting logic into the kernel.  The extensions would communicate so that the administrator could associate synchronization and kernel behavior with individual requests.

Extensions are fundamentally different from traditional software.  First, extensions require \emph{comprehensiveness}, or the ability to  interact with all of a system's software layers.  Second, extensions require \emph{dynamism}, or the ability to be added to a running system, so that the user can add monitoring or tune customization on the actual hardware and workloads that their system encounters.  Third, extensions require \emph{safety}, or the assurance that a buggy extension does not crash the system and that a subverted application cannot circumvent any extensions.  Safety is especially important for security monitoring extensions and for scenarios where the extension writer is not the original software developer of the extended system.  Finally, extensions require \emph{efficiency}, or the ability to execute at near native speed. 

Alas, existing system extension frameworks do not meet these needs.  Most systems are not comprehensive, as they only allow extensions to userspace applications (\eg dynamic binary instrumentation tools~\cite{Luk05, Nethercote07, Bruening12}, \texttt{LD\_PRELOAD}-based solutions, \etc) or only allow extensions to the operating system (\eg Nooks~\cite{Swift03}, VINO~\cite{Seltzer96}, SPIN~\cite{Bershad95}, linux kernel modules, \etc).   Extended Berkeley Packet Filters (eBPFs) are \emph{almost} comprehensive; they allow a user to extend the logic of the kernel and userspace applications, with the caveat that userspace extensions cannot modify userspace behavior. 
%\yusheng{behavior -> execution?}\arq{behavior just meaning "actions", but slightly fancier wording}
eBPF is dynamic, since users can extend a running system.  eBPF is safe since it uses a bytecode verifier to ensure that extensions cannot crash the system and executes extensions in the kernel to ensure that extensions are isolated from subverted applications.  However, eBPF is \emph{not} efficient, since its isolation approach requires multiple context switches to execute a userspace extensions.  For use cases with frequently-executed application extensions, such as the microservice monitoring example, this context switching overhead is substantial and can even lead to incorrect conclusions about the cause of performance degradation.

This paper presents \sys, a new system for comprehensive, dynamic, safe, and efficient extensions.  \sys{} is based upon two design principles.  First, the system maintains compatibility with the eBPF ecosystem so that it can build on eBPF's widespread adoption.  Second, the system resolves the inefficiency of the eBPF ecosystem by executing userspace extensions in the execution context of the extended application without requiring kernel involvement.  %Following these design principles, \sys{} enables the microservice monitoring use case to accurately measure performance and identify bottlenecks.

% and efficient full system extensions by allowing developers to ubiquitously use eBPF across their entire system.  \sys is compatible with the existing eBPF ecosystem, which makes it easy for developers to adopt.  Namely, developers extend kernel logic by hooking \texttt{kprobe}s, \texttt{kretprobe}s, and (\texttt{syscall}) \texttt{tracepoint}s; and extend application logic by hooking \texttt{uprobe}s and \texttt{uretprobe}s.  \sys provides a shared memory interface that allow extensions to communicate, whether the extensions are across system layers (\ie userspace extension that shares data with a kernel extension) or across processes (\ie a userspace extension from one process sharing data with a userpsace extension for another process).  \sys ensures extension safety---applications cannot interfere with extensions---and efficiency---extensions impose low runtime overhead.  Thus, with \sys, the microservice developer's monitoring is efficient and accurately measures performance.

% \yusheng{\sys provides a shared memory interface and data structs, so that extensions can communicate or aggregation data between different process or between kernel space and user space. For example, The kernel kprobe eBPF program, which runs in the kernel,  can update a variable in share memory to notify the  uprobe eBPF program, which runs in userspace the fsync operation is completed; The SSL/TLS unencrypted traffic data between different process can be obtained by uprobe, and analysis metrics without sending them all to a centric agent.} 
%\yusheng{I think it may not be accurate to describe bpftime as "API compatible"? The compatible of bpftime is more like compatible with syscalls. From the bpf program side, we provide compatible with helpers(like the syscall id) and maps(like stl lib in cpp); for the tracing application like bpftrace, we emulate bpf syscall behavior in linux systems.}


Realizing \sys{}'s design principles requires overcoming three challenges.  First, \sys{} requires an eBPF ecosystem to execute eBPF extensions in userspace; the ecosystem must be feature compatible with the existing eBPF.  \sys{} solves this challenge with an ``extra layer of indirection'', a software layer that interposes on all requests made to eBPF and forwards them to \sys{}.  Second, \sys{} requires trap-free attachment to an arbitrary program to ensure that userspace extensions can execute without kernel involvement.  \sys{} solves this challenge by using binary rewriting~\cite{Dinesh20, Chamith17, WilliamsKing20, fridagum} to inject userspace extensions into an application.  Finally, executing userspace extensions in the extended application threatens the safety of \sys{}, since it complicates the use of the existing in-kernel eBPF bytecode verifier and breaks the isolation between extensions and applications.  \sys{} solves this challenge by using a userspace eBPF bytecode verifier~\cite{Gershuni19} to prevent a buggy extension from harming an application, and by adopting work on intra-process hardware memory protections~\cite{Vahldiek19, xie2022cetis}, to prevent a subverted application from circumventing an extension. 

% The key to empowering \sys's safety and efficiency is a new technique called \tech. \tech enables efficient userspace extensions by executing them directly in the context of the original application, similar to binary instrumentation~\cite{Luk05}.  \tech's key innovation  is its low-overhead load-time approach to ensuring isolation between applications and extensions that contrasts with the expensive heavyweight sandbox approach employed in prior work.  Namely, \tech reuses existing eBPF verifiers~\cite{Gershuni19, linuxverifier} to ensure that each extension does not interfere with the program that it extends (\eg an extension cannot cause a program to have a null-pointer exception). \tech employs recent hardware features (\ie control-flow enforcement technology and memory protection keys) to ensure that the original application cannot interfere with the extension.

We implement \sys{} on linux.  The implementation requires \emph{no changes} to the operating system, since it is implemented entirely in the existing eBPF ecosystem.  Thus, a user can write extension against the well-known eBPF ecosystem, execute them on a regular linux kernel, and still benefit from \sys{}'s enhanced comprehensiveness and efficiency.

We illustrate the generality of \sys by using it in \numcases use-cases (\cref{sec:use_cases}).  The use-cases explore design tradeoffs, improve security, monitor performance, and improve system performance.  In particular, we use \sys to monitor a microservice application, create new durability configurations in Redis, cache meta data operations in FUSE, implement a ssl-supporting distributed tracing tool, monitor system calls for performance monitoring, and enhance Nginx's security.

We evaluate \sys's performance on the aforementioned \numcases use-cases.  We find that: \sys improves the throughput of profiling microservices by a factor of 1.5 compared to eBPF.  \sys{} enables a redis durability configuration that is 5 orders of magnitude more durable than the existing default setting, but only decreases throughput by ~10\%.  \sys enables FUSE caching that accelerates operations by orders of magnitude.  \sys{} reduces the overhead of SSL traffic monitoring by a factor of 1.67 compared to native eBPF.  And, \sys offers configurations that prevent monitoring overhead from affecting unmonitored processes, whereas eBPF is unable to do so.  We also use microbenchmarks to illustrate the key features that enable \sys's performance advantages and show \sys's compatibility with eBPF.

In sum, our contributions are:
\begin{smitemize}
\item An examination of system extension use-cases and the distillation of four design goals: comprehensiveness, dynamism, safety, and efficiency. 
\item \sys, a comprehensive, dynamic, safe, and efficient extension framework that maintains comparability with the eBPF ecosystem and adds efficiency by executing userspace extensions without kernel involvement.
\item An evaluation of a \sys showing that it is useful in \numcases use-cases and significantly more efficient than state-of-the-art alternatives.
\end{smitemize}

The rest of the paper is as follows.   We motivate system extensions and derive design goals (\cref{sec:motivation}).  We elaborate on the design principles and challenges in building the system (\cref{sec:challenges}).  Then, we discuss \sys's design (\cref{sec:design}) and implementation (\cref{sec:implementation}).  We discuss the \numcases use-cases (\cref{sec:use_cases}).   Next, we describe our evaluation (\cref{sec:evaluation}), related work (\cref{sec:related_work}), and conclude (\cref{sec:conclusion}).
